\bookchapter{第十六章：账号体系与云服务}

\chapterintro{我有个朋友，去年手机被偷了。他第一时间想的不是“手机多少钱”，而是“我里面的照片还能找回来吗”。很遗憾，他没开云备份。手机是身外之物，但账号里的数据，是真的丢了就再也回不来的东西。这章我们就来聊聊，怎么给你的数字生活上个真正的保险。}

\booksection{一、你的手机账号，到底管着多少东西？}

很多人觉得账号就是用来下载App的。其实，一个手机账号背后牵着的东西，比你想象的要多得多：

\begin{itemize}
  \item \textbf{通讯录}：你存了几年甚至十几年的电话号码。
  \item \textbf{照片和视频}：从孩子出生到毕业典礼，从第一次约会到结婚旅行。
  \item \textbf{备忘录}：你随手记的灵感、日记、账号密码（对，很多人把密码记在备忘录里）。
  \item \textbf{App购买记录}：花了几百块买的付费App和内购。
  \item \textbf{查找手机}：手机丢了之后，这是唯一能帮你定位、锁机、远程抹数据的工具。
\end{itemize}

如果你用的是苹果手机，以上所有东西都绑在你的Apple ID上。如果你用的是安卓，那就绑在你的华为/小米/OPPO/vivo账号上，还有你的谷歌账号。这些账号就是你数字生活的“户口本”，丢了户口本，很多东西就说不清了。

\booksection{二、三个必须做的事，现在就做}

\begin{enumerate}
  \item \textbf{开启双重认证}：不管你用的是苹果还是安卓，第一件事就是去账号设置里把"双重认证"或"两步验证"打开。这东西就像给你的门加了一把反锁——光有密码不够，还得你的手机收个验证码。就算有人在网上猜到了你的密码，没有你的手机，他也进不去。
  \item \textbf{绑定一个可靠的恢复方式}：双重认证很安全，但也有个风险——你自己的手机丢了怎么办？所以系统会提示你设置"恢复密钥"或"备用邮箱/手机号"。这个是你在紧急情况下拿回账号的唯一途径，最好记在纸上，别存在手机里。
  \item \textbf{确认"查找手机"已开启}：这个功能在手机丢了之后能帮你定位、锁定、播放声音、甚至远程擦除所有数据。是最后一道防线。很多人买手机的时候跳过了这一步，后面想用的时候发现没开，就晚了。
\end{enumerate}

\booksection{三、iCloud：苹果用户的“空气”}

用苹果手机的人，迟早要面对一个灵魂拷问：iCloud空间满了怎么办？

苹果免费给你5GB，说实话，这5GB连备份一次都不够。所以大部分人最后都会收到那个烦人的弹窗：“iCloud储存空间已满”。这时候你有三个选择：

\begin{itemize}
  \item \textbf{清理}：去设置 你的名字 iCloud 管理储存空间，看看是谁占了你的空间。一般来说，照片和备份是大头。可以先删掉不需要的旧备份。
  \item \textbf{升级}：50GB每个月6块钱，200GB每个月21块钱。个人觉得50GB对大多数人来说够了，主要是用来备份和存照片的缩略图。
  \item \textbf{曲线救国}：如果你不想花钱，可以把照片定期导出到电脑上，然后关掉iCloud照片同步。但这样的话，你就失去了照片自动备份的便利性。
\end{itemize}

一个重要的提醒：iCloud照片图库不是“网盘”——它是“同步”。你在手机上删了一张照片，iCloud上也会删。如果你想要的是“备份到云端然后清空手机”，那iCloud照片不是你要的东西，你需要的是iCloud云备份，或者第三方的网盘。

\begin{tcbinfo}{iCloud+：付费用户才能享受的隐私三件套}
只要升级到任意付费iCloud档位（50GB/200GB/2TB等），就自动成为\textbf{iCloud+}用户，解锁三个隐私保护功能：

\begin{itemize}
  \item \textbf{Private Relay（私密中继）}：苹果版的"VPN"。Safari浏览的流量经苹果服务器中转后再发往网站，没有人（包括苹果）能同时看到你是谁、你访问了什么。适合在公共Wi-Fi下保护隐私。设置路径：\texttt{设置 > 你的名字 > iCloud > 私密中继}。
  \item \textbf{Hide My Email（隐藏邮件）}：在网站/App注册时生成一次性随机邮箱地址，邮件会自动转发到你的真实邮箱。不想再收这家网站的邮件，直接关掉这个随机地址就行，真实邮箱不暴露。设置路径：\texttt{设置 > 你的名字 > iCloud > 隐藏邮件}。
  \item \textbf{自定义电子邮件域名}：可以用你自己的域名（比如@你的姓.com）收发iCloud邮件，适合个人品牌或家庭邮箱管理。
\end{itemize}

提示：Private Relay在某些国家/地区不可用，国内用户可能需要其他方案补充。

\end{tcbinfo}

\begin{tcbwarn}{iCloud高级数据保护（ADP）——真正的端到端加密}
2022年底苹果推出了\textbf{iCloud高级数据保护（Advanced Data Protection，ADP）}，把端到端加密（E2EE）从原来的健康、密码等少数数据，扩展到了\textbf{绝大部分iCloud数据}——照片、备忘录、iCloud云备份、设备信息等共23类数据。

\begin{itemize}
  \item \textbf{什么是E2EE}：加密密钥只存在你的设备上，连苹果服务器都没有。就算苹果被黑客攻破，或者执法部门要求调取，苹果也拿不出能解密的明文数据。
  \item \textbf{开启路径}：\texttt{设置 > 你的名字 > iCloud > 高级数据保护}，按提示设置恢复联系人/恢复密钥（\textbf{必做}，丢了密钥你的数据也救不回来）。
  \item \textbf{注意事项}：(1)开启后务必妥善保管恢复密钥，否则账号丢失时无法找回数据；(2)共享相簿、iCloud共享图库等少数功能不参与E2EE；(3)中国大陆区Apple ID暂不支持ADP，需要切换到其他地区账号。
\end{itemize}

一句话建议：如果你的账号里有重要隐私数据，且账号地区支持，强烈建议开启ADP——这是普通人能获得的最强云端数据保护。

\end{tcbwarn}

\booksection{四、同步和备份，别搞混了}

这两个词听起来差不多，但后果完全不同：

\begin{itemize}
  \item \textbf{同步}是双向的：你在手机上删了，云端也删了；你在云端删了，手机上也跟着删。iCloud照片、通讯录、备忘录这些都是同步。好处是多设备之间保持一致，坏处是"删了就真没了"。
  \item \textbf{备份}是单向的：手机把当前状态拍一张"快照"存到云端。之后你在手机上删了东西，云端备份里的那份还在。iCloud云备份、电脑本地备份，都是备份。好处是数据更安全，坏处是占空间，而且不能实时更新。
\end{itemize}

\booksection{五、安卓用户的账号矩阵}

安卓这边的情况比苹果复杂一些，因为你的数据可能分散在好几个账号里：

\begin{itemize}
  \item \textbf{品牌账号}（华为/小米/OPPO/vivo）：管着你的云空间、查找手机、应用商店、钱包。每个品牌都有自己的云服务，通常是基础免费+扩容收费的模式。
  \item \textbf{谷歌账号}：如果你用Gmail、Google相册、Google通讯录，这些东西都在谷歌账号里。谷歌免费给15GB，比苹果大方，但也是和Gmail、Drive共用的。
  \item \textbf{微信/QQ}：它们是独立于手机系统的。你的聊天记录和微信收藏，既不在品牌账号里也不在谷歌账号里，只存在微信自己的服务器上。
\end{itemize}

所以安卓用户做备份的时候，要分别检查：品牌云备份、谷歌备份、微信聊天记录备份，三者缺一不可。

\booksection{六、家庭共享：让家人也受益}

苹果的“家人共享”可以让你和家人共享App购买、音乐、电影、iCloud空间，还能共享位置。最多可以加6个人。设置方法是：设置 你的名字 家人共享。华为和小米也有类似的功能，在账号设置里能找到“家庭共享”或“家庭云”。

一个很实用的场景：给爸妈的手机上设置好家人共享，这样你就能随时看到他们的位置（前提是他们同意），万一老人走丢了或者手机联系不上了，能有个定位。

\begin{tcbwarn}{账号被盗的后果比你想象的要严重}
如果黑客拿到了你的账号密码，他可以：远程锁定你的手机并勒索赎金、查看你的照片和通讯录、用你的支付方式买东西。更麻烦的是，如果你的账号绑定了太多东西，找回账号本身就是一场噩梦。所以，双重认证，现在就去开。

\end{tcbwarn}

\begin{tcbpractice}{实操：给你的账号做一次全面体检}
\begin{enumerate}
  \item 打开手机设置，进入账号页面，确认双重认证已开启。
  \item 检查"受信任的电话号码"和"恢复密钥"是否已设置，密钥记在纸上。
  \item 进入iCloud或云空间管理，看看哪些数据在同步，关闭不需要的。
  \item 确认"查找我的手机"已开启。
  \item 如果你用的是安卓，检查品牌云备份、谷歌备份、微信备份三者是否都覆盖到了。
\end{enumerate}

\end{tcbpractice}
